\section{Evaluation}
\label{sec:evaluation}

\textbf{Implementation.} We implemented the presented layered BRDF for the isotropic and the anisotropic case in GLSL fragment shaders. Therefore, we wrote our own C++ OpenGL real-time rendering engine. Some features of our engine are:
\begin{itemize}
    \item obj model loading
    \item saving/loading scenes and material presets
    \item easy switching between shaders and models
    \item loading cubemaps
\end{itemize}
Our implementation of the layered BRDF does not rely on precomputation, instead we just use the Fresnel term as an approximation of FGD. Furthermore, we limited the number of layers to 4 (could be increased of course) and did not implement the doubling algorithm, which is not really a limitation. For each layer $\eta$, $\kappa$ and $\alpha$ can also be textured. The anisotropic shader is using a two-component vector for the variance. In our engine are currently two modes of lighting possible: a directional light or image based lighting (IBL). For IBL in the layered shaders we are using Monte Carlo integration with importance sampling as described in Section \ref{sec:anisotropic_interfaces}. The hammersley sequence is used for generating low-discrepancy points in order to sample the VNDF.

%

Each layer is parametrized by a three-component complex index of refraction $\eta + i \kappa$, a scalar/two-component roughness $\alpha$, a scalar depth $h$, a three-component absorption cross-section $\sigma_a$, a three-component scattering cross-section $\sigma_s$ and a phase function anisotropy factor $g$.\\

%

\noindent\textbf{Results.} A lot of the desired goals mentioned in Section \ref{sec:introduction} are actually fulfilled by the presented model and our implementation of it:
\begin{itemize}
    \item the model is physically based because of using the microfacet model following the GGX distribution and because it tries to approximate the physical light transport
    \item it supports an arbitrary number of layers
    \item multiple scattering is accounted
    \item if incorporated in FGD, the model is energy conserving
    \item texturing and importance sampling is supported
    \item there is no or only little precomputation and storage needed
    \item the model can be evaluated in real-time
\end{itemize}

However, there are some limitations of the general accuracy and regarding the effects the model is able to reproduce. Still, the model is able to convincingly reproduce a variety of complex layered materials. For instance metal foil, metallic paint, frosted metal, clear coating and gold coating, see Figure \ref{fig:possible_materials}. 

\begin{figure}[!h]
    \includegraphics[width = \columnwidth, keepaspectratio]{christmas_balls.png}
    \caption{Textured layers.}
    \label{fig:textured_materials}
\end{figure}

The scene shown in Figure \ref{fig:textured_materials} illustrates how the model can be used to represent complex materials using textures for the index of refraction and roughness. Furthermore, Figure \ref{fig:match} shows a comparison  between our implementation, the offline implementation of \textcite{belcour2018} and the reference Belcour used. One can see that, apart from missing energy conservation resulting in a darker appearance and slight differences in the environment map, our implementation closely matches both of them. Effects of absorption in participating media can be seen in Figure \ref{fig:absorption}.

\begin{figure}[!h]
    \centering
    \begin{subfigure}[t]{0.3\columnwidth}
        \centering
        \includegraphics[width = \columnwidth, keepaspectratio]{belcour_our_try_match.png}
        \caption{Our implementation}
    \end{subfigure}
    \begin{subfigure}[t]{0.3\columnwidth}
        \centering
        \includegraphics[width = \columnwidth, keepaspectratio]{fromBelcour/belcour_offline.png}
        \caption{Offline model of \textcite{belcour2018}}
    \end{subfigure}
    \begin{subfigure}[t]{0.3\columnwidth}
        \centering
        \includegraphics[width = \columnwidth, keepaspectratio]{fromBelcour/belcour_reference.png}
        \caption{Reference \textcite{belcour2018} used}
    \end{subfigure}
    \caption{Our implementation results in a darker appearance but still closely matches the implementation of \textcite{belcour2018} and the reference they used. We used the Images of Figure 17 from \textcite{belcour2018}.}
    \label{fig:match}
\end{figure}

\begin{figure}[!h]
    \includegraphics[width = \columnwidth, keepaspectratio]{absorption_effects.png}
    \caption{Effects of absorption in participating media, leading to angle dependent saturation, can be reproduced with our implementation.}
    \label{fig:absorption}
\end{figure}

%

In Figure \ref{fig:white_furnace} we show a white furnace test. It demonstrates that because of just using the Fresnel term as an approximation for FGD instead of incorporating some sort of method to account for multiple scattering and because the GGX BRDF itself is not energy conserving, our implementation is not energy conserving.

\begin{figure}[!h]
    \centering
    \begin{subfigure}[t]{0.47\columnwidth}
        \centering
        
\includegraphics[width = \columnwidth, keepaspectratio]{white_furnace_a.png}
        \caption{$\alpha_0=0.001$}
    \end{subfigure}
    \begin{subfigure}[t]{0.47\columnwidth}
        \centering
        
\includegraphics[width = \columnwidth, keepaspectratio]{white_furnace_b.png}
        \caption{$\alpha_0=1.0$}
    \end{subfigure}
    \caption{White-furnace test for a three layer configuration: $\eta_0 = 1.3$, $\eta_1=1.1$, $\eta_2=0.0$, $\kappa_2 = 1.0$, $\alpha_1=0.001$, $\alpha_2=0.001$ }
    \label{fig:white_furnace}
\end{figure}

%

\textit{Performance.} All results are computed on an 8 core Intel i7-1165G7 with 16 GB of RAM and an integrated Mesa Intel Xe Graphics GPU measured using OpenGL timer queries. Render times compared to the standard GGX shader are reported in Figure \ref{fig:performance_ggx_comparison}. It can be seen, that our layered shader is on average approximately 2.9 times slower, however it converges to approximately 10ms per frame for $1920\times1080$, thus it is fast enough for real-time usage.

\begin{figure}[!h]
    \includesvg[width = \columnwidth, keepaspectratio]{performance_plot.svg}
    \caption{Render times of a full-screen plane with our implementation of the layered shader with two layers compared to the standard GGX shader. Evaluated for different window sizes.}
    \label{fig:performance_ggx_comparison}
\end{figure}

Figure \ref{fig:performance_layers} shows, as expected, that the runtime is linear in the number of layers, as during evaluating the adding-doubling algorithm we iterate over all layers. Our implementation was written with readability in mind instead of performance optimization. Therefore, the render times could probably be decreased.

\begin{figure}[!h]
    \includesvg[width = \columnwidth, keepaspectratio]{performance_layer_plot.svg}
    \caption{Render times of a full-screen plane with our implementation of the layered shader for different numbers of layers at $1920\times1007$ window size.}
    \label{fig:performance_layers}
\end{figure}

%

In Figure \ref{fig:performance_ibl} a comparison of sample counts and render times is displayed for image based lighting using Monte Carlo integration of the anisotropic layered BRDF. This shows the limited usability of the anisotropic shader for real-time applications, as the needed sample counts for visually pleasing results are too high. The resulting images can be seen in Appendix \ref{sec:appendix_b}. 

\begin{figure}[!h]
    \includesvg[width = \columnwidth, keepaspectratio]{performance_ibl_plot.svg}
    \caption{Render times of a sphere with our implementation of the anisotropic layered shader using image based lighting. Evaluated for different numbers of samples with different roughness values at $800\times800$ window size. The material consists of two layers, the roughness of the top layer is changed, the roughness for the bottom layer is $\alpha_1=(0.001,0.001)$}.
    \label{fig:performance_ibl}
\end{figure}

\begin{figure*}
    \centering
    \begin{subfigure}[t]{0.16\textwidth}
        \centering
        \includegraphics[width = \textwidth, keepaspectratio]{metal_foil.png}
        \caption{Metal foil}
    \end{subfigure}
    \begin{subfigure}[t]{0.16\textwidth}
        \centering
        \includegraphics[width = \textwidth, keepaspectratio]{metallic_paint.png}
        \caption{Metallic paint}
    \end{subfigure}
    \begin{subfigure}[t]{0.16\textwidth}
        \centering
        \includegraphics[width = \textwidth, keepaspectratio]{frosted_metal.png}
        \caption{Frosted metal}
    \end{subfigure}
    \begin{subfigure}[t]{0.16\textwidth}
        \centering
        \includegraphics[width = \textwidth, keepaspectratio]{rough_on_rough.png}
        \caption{Rough on rough}
    \end{subfigure}
    \begin{subfigure}[t]{0.16\textwidth}
        \centering
        \includegraphics[width = \textwidth, keepaspectratio]{clear_coat.png}
        \caption{Clear coated}
    \end{subfigure}
    \begin{subfigure}[t]{0.16\textwidth}
        \centering
        \includegraphics[width = \textwidth, keepaspectratio]{gold.png}
        \caption{Gold coated}
    \end{subfigure}
    \caption{We are able to reproduce a variety of appearances with our implementation. All of them consist of two layers, a dielectric on top of a conductor, with varying roughnesses.}
    \label{fig:possible_materials}
\end{figure*}
